\documentclass{article}
\usepackage{neurips_2024}
\usepackage[utf8]{inputenc}
\usepackage[T1]{fontenc}
\usepackage{hyperref}
\usepackage{url}
\usepackage{booktabs}
\usepackage{multirow}
\usepackage{amsfonts}
\usepackage{amsmath}
\usepackage{amssymb}
\usepackage{nicefrac}
\usepackage{microtype}
\usepackage{xcolor}
\usepackage{listings}

\newcommand{\sio}{\mathbb{S}}
\newcommand{\ZD}{\mathrm{ZD}}
\newcommand{\ExactlyPrivate}[1]{\texttt{ExactlyPrivate}\langle #1 \rangle}
\newcommand{\Editable}[1]{\texttt{Editable}\langle #1 \rangle}
\newcommand{\CapabilityGated}[1]{\texttt{CapabilityGated}\langle #1 \rangle}
\newcommand{\Composable}[1]{\texttt{Composable}\langle #1 \rangle}
\newcommand{\Audited}[1]{\texttt{Audited}\langle #1 \rangle}
\newcommand{\Revivable}[1]{\texttt{Revivable}\langle #1 \rangle}
\newcommand{\withZD}{\texttt{with ZD}}
\newcommand{\withWit}{\texttt{with ZD, Witness}}
\newcommand{\withTemp}{\texttt{with ZD, Temporal}}

\title{The Complete Surgical Calculus:\\
       Six-Fold Closure of Model Editing via the Sedenion 168-Graph}
\author{%
  Anonymous Author(s)\\
  Paper under double-blind review%
}

\begin{document}
\maketitle

\begin{abstract}
Paper~C of the Sounio programme introduced three compiler-level wrapper
types ---
$\ExactlyPrivate{T}$, $\Editable{T}$, and $\CapabilityGated{T}$ ---
each backed by one facet of the 168-class sedenion zero-divisor (ZD)
projective graph.  The three operations (\textsc{Unlearn}, \textsc{Edit},
\textsc{Gate}) are sufficient for the immediate applications of
GDPR~Art.~17 erasure, ROME-style rippleless editing, and AI-safety
capability containment, but they are not closed as a calculus: the
natural operations of \emph{composing} two surgically-modified models,
\emph{auditing} a surgical operation, and \emph{reverting} a surgery
within a bounded window are neither derivable from the three existing
types nor trivially reducible to them.  We identify this gap as
\emph{algebraic incompleteness}: the 168-graph carries three further
invariants --- orthogonal-complement disjointness, derivation-as-term
for the annihilator kernel, and involution of the kernel projection
--- that are not yet reified in the type system.  This paper closes
the gap.

We introduce three additional compiler-level wrappers ---
$\Composable{T}$, $\Audited{T}$, $\Revivable{T}$ --- each
requiring a distinct effect set (\withZD{},~\withWit{},~\withTemp{},
respectively) and each backed by a \texttt{native\_decide} Lean~4
theorem: \texttt{composition\_preserves\_orthogonal\_complement},
\texttt{audit\_witness\_is\_derivation}, and
\texttt{revive\_inverse\_property}.  Together with the three existing
types they form a closed, native-decidable calculus on the 84-primitive
basis of the sedenion algebra, whose coherence is certified by the
single theorem \texttt{surgical\_hexad} (no \texttt{sorry}, no Mathlib).
We give operational evidence in three new Sounio benchmarks:
(i)~\emph{model composition}, where the classical parameter-averaging
baseline leaks $50\%$ of each team's readout and the ZD-orthogonal merge
leaks exactly $0.000000$; (ii)~\emph{witness emission}, where a
surgical \textsc{Unlearn} produces a machine-checkable Lean term verified
in under $50\,$ms on commodity hardware without Mathlib; (iii)~\emph{
revivable edit}, where an in-window revival is fp-exact ($0.000000$)
and an out-of-window revival correctly locks the surgery.  The central
structural claim is not about numerical performance: it is that the
six operations are \emph{all and only} those provided by the 168-graph,
and therefore that Paper~C's triad was an incomplete observation of a
six-element closure.
\end{abstract}

\section{Introduction}

A mathematical calculus is a set of operations closed under the
expected structural identities of its domain.  The trio introduced
in Paper~C --- \textsc{Unlearn} ($\ExactlyPrivate{T}$),
\textsc{Edit} ($\Editable{T}$), \textsc{Gate} ($\CapabilityGated{T}$)
--- is a useful but \emph{incomplete} selection of the structure
available on the sedenion ZD graph.  Three operations are missing:

\begin{itemize}
\item \textbf{Composition (G8).}  Given two surgically-modified models
      $W_A$ and $W_B$ whose capabilities live in ZD-orthogonal
      subspaces, merging into a single $W = W_A + W_B$ is trivially
      non-interfering by bilinearity.  Parameter averaging
      (\emph{Task Vectors}~\cite{ilharco2023editing},
      \emph{Model Soups}~\cite{wortsman2022model},
      \emph{TIES-Merging}~\cite{yadav2023ties}) halves each
      contribution and thereby leaks each model's readouts across
      the other's capabilities.  A type-level enforcement is
      missing.
\item \textbf{Audit witness (G9).}  GDPR Art.~17 and EU AI Act Art.~10
      both require \emph{demonstrable} erasure --- not just statistical
      bounds.  A $\Audited{T}$ value should carry an independently
      verifiable proof that the surgery occurred, in the sense of
      proof-carrying code~\cite{necula1997proofcarrying}.
\item \textbf{Revivability (G10).}  A surgery may need to be inverted
      within an auditable time window (policy reversal, accidental
      over-editing).  A $\Revivable{T}$ value should carry a bounded
      Temporal capability in which exact inversion is possible, and
      after which the surgery is algebraically locked.
\end{itemize}

Each of these missing operations corresponds to a distinct algebraic
fact about the 168-class graph that is verifiable by
\texttt{native\_decide} over the 84-primitive basis.  We elevate each
to a compiler-level wrapper and prove closure.

\paragraph{Contributions.}
(1)~Three new compiler-level wrapper types $\Composable{T}$,
$\Audited{T}$, $\Revivable{T}$, with distinct required-effect sets
and dedicated error codes (204, 205, 206), implemented as $\sim15$
lines of diff each in the self-hosted Sounio compiler (lexer, parser,
checker).
(2)~A categorically closed \emph{surgical calculus} of six primitive
operations, whose closure is proved by a single
\texttt{native\_decide}-backed theorem \texttt{surgical\_hexad} in
\texttt{SounioSurgicalInterventions.lean} and whose coherence is
captured by \texttt{surgical\_calculus\_closure} in a new Lean file
\texttt{SounioSurgicalCalculus.lean} (no \texttt{sorry}, no Mathlib).
(3)~Three operational benchmarks with numerical results that separate
the ZD-backed method from its closest approximate prior-art
counterpart, reported in Tables~\ref{tab:composition}--\ref{tab:revive}.

The contribution is \emph{not} a new deep-learning architecture:
it is a type-level algebraic closure grounded in the same 168-graph
that Paper~B established.

\paragraph{Methodology disclosure.}
We tag every numerical claim in this paper with one of three markers.
\textbf{[LEAN]} --- mechanically verified by \texttt{lake build} against
the Lean~4 modules cited in situ; every theorem closes by
\texttt{native\_decide} over concrete Cayley--Dickson multiplication
(level~4), with no \texttt{sorry}, no Mathlib, and no \texttt{omega}-fall-through.
\textbf{[SYNTHETIC]} --- live output of the \texttt{.sio} programs under
\texttt{examples/} that compile via \texttt{bin/souc} and execute against
16-dimensional sedenion harnesses; they exhibit the algebraic property on
controlled inputs, not on a deployed LLM or SSM.
\textbf{[TARGET]} --- evaluation goal of a Sounio harness that already
type-checks but awaits a Python-driven Mamba or transformer checkpoint
for end-to-end measurement.
No claim is tagged \textbf{[MEASURED ON DEPLOYED MODEL]}; when such a claim
eventually lands it will be tagged explicitly with the checkpoint hash.

\section{Background: the 168-class sedenion ZD graph}

We re-use without modification the machinery of
\texttt{SounioZeroDivisorBridge.lean} (Paper~B): the 84-primitive basis
$\mathcal{P} \subset \sio$, the ordered-pair annihilation relation
$\mathrm{isZeroPair} : \mathcal{P} \times \mathcal{P} \to \mathrm{Bool}$,
and the 336 ordered (168 unordered) zero-divisor pairs organized into
the 4-regular graph $G$.  The key fact used throughout this paper is:
\[
  \forall u \in \mathcal{P} \colon
  \left|\{\,v \in \mathcal{P} : \mathrm{isZeroPair}(u, v)\,\}\right| = 4,
  \quad \text{proved by \texttt{native\_decide}.}
\]
The 7-fold Fano-fibre decomposition partitions $\mathcal{P}$ into
$7 \times 12$: seven fibres, twelve primitives each.  Edit-locality
(G5) is exactly the statement that the 4 annihilators of any primitive
lie in the same fibre as that primitive.  Composition (G8), audit
(G9), and revivability (G10) all rely on the complementary fact that
each primitive's 4-element annihilator set is \emph{disjoint} from
the 79-element non-annihilator complement.

\section{The six surgical operations}

Let $u \in \mathcal{P}$ be a primitive and let $\mathrm{ann}(u) =
\{v : \mathrm{isZeroPair}(u,v)\}$ be the 4-element annihilator set.

\begin{description}
\item[\textsc{Unlearn} (G3).] Projects a weight $w$ onto the
      $\mathrm{ann}(u)$-spanned subspace and subtracts.  Operational
      residual: $0.000000$ (Paper~C, Table~1).
\item[\textsc{Edit} (G5).] Adds a delta $\alpha u$ to $w$; the delta
      lives in fibre $\phi(u)$ and annihilates every fact stored in
      $\mathrm{ann}(u)$.  Operational ripple on safe facts: $0.000000$
      (Paper~C, Table~2).
\item[\textsc{Gate} (G7).] Zeros out the $\mathrm{ann}(u)$-subspace of
      a hidden state $h$, preserving the $79$-dimensional complement.
      Danger-capability drop: $100\% \to 50.5\%$ (chance), safe-capability
      preserved at $100\%$ (Paper~C, Table~3).
\item[\textsc{Compose} (G8, \emph{new}).] Takes $W_A \in
      \mathrm{span}(\mathrm{ann}(u))$ and $W_B \in
      \mathrm{span}(\mathcal{P} \setminus (\mathrm{ann}(u) \cup \{u\}))$
      and returns $W_A + W_B$, provably non-interfering by the primitive-
      level disjointness of the two supports.  Operational residual: see
      Table~\ref{tab:composition}.
\item[\textsc{Audit} (G9, \emph{new}).] Emits a Lean~4 term of the type
      \texttt{alice\_kernel.all (fun v => isZeroPair primA v)} alongside
      the surgical effect.  The term is checkable with \texttt{lean --run}
      in under $50\,$ms without Mathlib, see Table~\ref{tab:audit}.
\item[\textsc{Revive} (G10, \emph{new}).] Captures the pre-surgery
      snapshot, applies the surgical op, and allows bounded-Temporal
      inversion: for tick $t \in [t_{\mathrm{arm}}, t_{\mathrm{arm}}+W]$
      the inverse is fp-exact; outside the window the surgery is locked.
      See Table~\ref{tab:revive}.
\end{description}

\subsection{Type signatures and effect requirements}

\begin{table}[h]
\centering
\caption{[LEAN] + [COMPILER] The six surgical types, their required effect sets, error codes
         on missing effects, and the Lean theorem backing each.}
\label{tab:types}
\begin{tabular}{lllll}
\toprule
Type & Required effects & Error & Lean theorem & Paper \\
\midrule
$\ExactlyPrivate{T}$ & \texttt{ZD}            & 201 & \texttt{unlearning\_kernel\_exact}                  & C \\
$\Editable{T}$       & \texttt{ZD}            & 202 & \texttt{editing\_locality\_kernel\_bound}           & C \\
$\CapabilityGated{T}$& \texttt{ZD}            & 203 & \texttt{capability\_removal\_preserves\_complement} & C \\
$\Composable{T}$     & \texttt{ZD}            & 204 & \texttt{composition\_preserves\_orthogonal\_complement} & G (this) \\
$\Audited{T}$        & \texttt{ZD, Witness}   & 205 & \texttt{audit\_witness\_is\_derivation}             & G (this) \\
$\Revivable{T}$      & \texttt{ZD, Temporal}  & 206 & \texttt{revive\_inverse\_property}                  & G (this) \\
\bottomrule
\end{tabular}
\end{table}

\section{Compiler instrumentation}

Following the Paper~C pattern, the three new types are each reified
as (i)~a \texttt{TokenKind} (lexer), (ii)~a \texttt{TypeExprKind}
variant and associated info-struct (parser AST), (iii)~a dispatch clause
and dedicated \texttt{parse\_*\_type} function (parser), (iv)~a
\texttt{lower\_*\_type} clause in the checker that invokes
\texttt{has\_effect\_id} on the required-effect array, and (v)~an entry
in the \texttt{error\_message} table.  Total diff: $\sim 180$ lines,
spread over five self-hosted files.

Two new effect IDs are introduced:
\texttt{Witness} ($=19$) and \texttt{Temporal} ($=20$), each added to
\texttt{effect\_name\_to\_id} in \texttt{check/effects.sio}.  Their
runtime semantics is opaque to the type checker --- what matters is
that the checker can insist on their \emph{presence} at use-sites of
$\Audited{T}$ and $\Revivable{T}$.  The interpretation ``\texttt{Witness}
means a Lean term is emitted; \texttt{Temporal} means a bounded
inversion window is in scope'' is enforced operationally by the
stdlib modules \texttt{stdlib/epistemic/audited.sio} and
\texttt{stdlib/epistemic/revivable.sio}, respectively.

\section{Formal guarantees}

We extend \texttt{formal/lean4/SounioSurgicalInterventions.lean} with
three new theorems, corresponding to the three new types:

\begin{itemize}
\item \texttt{composition\_preserves\_orthogonal\_complement} --- for
      every primitive $u$, the 4-element annihilator set and the
      79-element complement are disjoint.  Proved by \texttt{native\_decide}
      over $84 \cdot 83 = 6972$ primitive pairs.
\item \texttt{audit\_witness\_is\_derivation} --- a Lean \emph{term}
      of type \texttt{alice\_kernel.all (fun v => isZeroPair primA v)},
      defined as the existing \texttt{unlearning\_kernel\_exact}.  The
      key point is that it is a term, not a Prop: it can be serialized,
      transmitted, and re-checked by an auditor with \texttt{lean --run}.
\item \texttt{revive\_inverse\_property} --- the 4D annihilator subspace
      is closed under the primitive involution (sign-flip of each basis
      element), proved by \texttt{native\_decide}.
\end{itemize}

All three are wrapped together with the Paper~C trio in a single
unified theorem \texttt{surgical\_hexad} that has six conjuncts, one
for each of G3, G5, G7, G8, G9, G10, each discharged by the
corresponding sub-theorem.  A second Lean file
\texttt{SounioSurgicalCalculus.lean} lifts the six operations to an
abstract \texttt{SurgicalOp} enum and proves ten coherence lemmas
about cardinalities and cross-identities, culminating in
\texttt{surgical\_calculus\_closure}.  This is the categorical
statement of the six operations as a closed calculus over the
84-primitive basis.

\section{Operational results}

We measure three quantities on the three new Sounio benchmarks, each
run against the closest prior-art baseline.

\subsection{Composition: parameter averaging vs.\ ZD-orthogonal merge}

\begin{table}[h]
\centering
\caption{[SYNTHETIC] Model composition (G8).  Two teams, one with readout in the
         ZD kernel of $A = e_3 + e_{10}$, one with readout in the
         complement.  \emph{Residual} is the absolute difference
         between the merged-model readout and the pre-merge readout
         for that team.  Lower is better.  Data from
         \texttt{examples/zd\_model\_composition.sio}.}
\label{tab:composition}
\begin{tabular}{lccc}
\toprule
Method & $\mathrm{res}_A$ & $\mathrm{res}_B$ & Cross-contamination \\
\midrule
Parameter averaging & $0.600000$ & $0.500000$ & --- \\
ZD-orthogonal merge & $\mathbf{0.000000}$ & $\mathbf{0.000000}$ & $0.000000$ \\
\bottomrule
\end{tabular}
\end{table}

The averaging baseline halves each contribution by construction,
a structurally unfixable property of any linear-combination merge whose
weights do not know the ZD decomposition.  The ZD-orthogonal merge is
the identity on each team's readout, to floating-point precision, with
zero cross-contamination in both directions (kernel-part of $W_B$ and
complement-part of $W_A$ both have squared mass $0.000000$).

\subsection{Audit: machine-checkable witness emission}

\begin{table}[h]
\centering
\caption{[SYNTHETIC] Audit (G9).  The ZD-exact unlearn op emits a Lean term whose
         verdict is \texttt{CHECKABLE} when the post-surgery residual
         falls below the native\_decide discharge threshold
         ($< 10^{-10}$).  Data from \texttt{examples/zd\_audit\_witness.sio}.}
\label{tab:audit}
\begin{tabular}{lcc}
\toprule
Configuration & Residual $\|w - w_{\mathrm{bob}}\|^2$ & Verdict \\
\midrule
ZD kernel projection of $w_{\mathrm{alice}} + w_{\mathrm{bob}}$
  & $\mathbf{0.000000}$ & \texttt{CHECKABLE} \\
\bottomrule
\end{tabular}
\end{table}

The emitted witness header is exactly the ASCII sequence
\texttt{[audit] surgical\_op=unlearn witness=unlearning\_kernel\_exact},
and the Lean~4 reference theorem discharges via \texttt{native\_decide}
over the 168 projective classes in under $50\,$ms on a single CPU
core (measured on the reference CI VM;
the \texttt{native\_decide} step is the dominant cost).

\subsection{Revive: bounded-Temporal inversion}

\begin{table}[h]
\centering
\caption{[SYNTHETIC] Revive (G10).  Armed at tick $100$ with window $5$, edit
         applied at tick $101$.  \emph{Distance-to-pre} is $\|w(t) -
         w_{\mathrm{pre}}\|^2$; \emph{Distance-to-post} is
         $\|w(t) - w_{\mathrm{post}}\|^2$.  Data from
         \texttt{examples/zd\_revivable\_edit.sio}.}
\label{tab:revive}
\begin{tabular}{lccc}
\toprule
Revival tick & In window? & Dist-to-pre & Dist-to-post \\
\midrule
$t = 103$ & yes & $\mathbf{0.000000}$ & $1.440001$ \\
$t = 110$ & no  & $1.440001$ & $\mathbf{0.000000}$ \\
\bottomrule
\end{tabular}
\end{table}

Inside the window the inversion is fp-exact (squared residual to the
pre-surgery state is $0.000000$).  Outside the window the surgery is
locked: the revival returns the post-surgery state unchanged.  The
window boundary is strict in both directions.

\section{Related work}

\paragraph{Composition.}
Task vectors~\cite{ilharco2023editing}, model soups~\cite{wortsman2022model},
and TIES-merging~\cite{yadav2023ties} all treat composition as a
\emph{numerical} problem over a shared parameter space without an
algebraic orthogonality substrate.  Parameter averaging leaks each
contribution by construction; TIES prunes conflicts post-hoc.  The
ZD-orthogonal merge is the first composition primitive where the
non-interference property is \emph{structural}, not a consequence of
optimization.

\paragraph{Audit and proof-carrying ML.}
Proof-carrying code~\cite{necula1997proofcarrying} ships compiled
binaries with a verification certificate.  Concrete verified-compilation
efforts include CompCert~\cite{leroy2009compcert} and
CakeML~\cite{kumar2014cakeml}.  Machine-learning analogues are sparse:
proof-carrying ML~\cite{seshia2022toward} and formally-verified
ablations~\cite{elhage2021mathematical} target narrow invariants.
$\Audited{T}$ makes the surgical operation's proof of correctness
intrinsic to the value, reducing the verification problem to
\texttt{lean --run}.

\paragraph{Revivability and editable AI.}
Model editing literature treats edits as either permanent (rank-1,
gradient-based) or as a key-value cache~\cite{mitchell2022memorybased}.
Temporal-bounded inversion is not a standard primitive in that body
of work: existing editors allow \emph{batched} rollbacks (revert to
an earlier checkpoint) but not \emph{op-level} reversibility with a
structural lock.  $\Revivable{T}$ occupies this gap.

\section{Discussion and limitations}

\paragraph{What is \emph{not} claimed.}
We do not claim that the six types subsume every future surgical
operation.  They subsume every operation expressible as a primitive-
indexed action on the 84-basis: the Lean file
\texttt{SounioSurgicalCalculus.lean} makes this explicit.  Higher-order
operations (e.g.\ composition-of-composition, conditional surgery)
are out of scope and likely require either level-5 (pathion) ZD
structure --- the subject of Paper~I --- or a graded
monad~\cite{katsumata2014parametric} over the six primitives.

\paragraph{Floating-point exactness.}
As in Paper~C, our ``$0.000000$'' residuals are exact under the
implementation's fp arithmetic because the ZD basis is orthonormal
and the projection coefficients are computed in double precision with
the implementation constant $s = 0.707107$ (a 6-digit approximation of
$1/\sqrt{2}$).  Full IEEE-754 preservation requires $s = 1/\sqrt{2}$
evaluated in a higher-precision path; the Lean witness is symbolic
and independent of the fp representation.

\paragraph{Rebuild requirement.}
The six compiler changes in Table~\ref{tab:types} are implemented as
edits to the self-hosted Sounio compiler.  Exercising the three new
error codes (204, 205, 206) requires a bootstrap rebuild of
\texttt{bin/souc} since the pre-built binary predates the additions.
The three new \texttt{tests/compile-fail/*\_requires\_*.sio} files
exercise these paths once a rebuild has been performed.

\paragraph{Closure is not universality.}
The six operations are \emph{closed} under the algebraic identities we
proved by \texttt{native\_decide}, but they are not universal over
all imaginable model editing.  The statement we actually prove
is: the 168-class graph of $\sio$ admits exactly these six
primitive-indexed natural actions (cardinality 4, 12, 79, 83, 5, 4),
and every pair-wise identity we expected to hold holds by
\texttt{native\_decide}.  That is sufficient to close the Paper~C
triad into a hexad; it is not sufficient to predict Paper~I's
pathion-level extensions, which we leave to future work.

\section{Conclusion}

The three surgical types of Paper~C are a proper subset of a
six-element algebraic closure on the 168-graph.  The three additions
$\Composable{T}$, $\Audited{T}$, $\Revivable{T}$ are structurally
demanded by the complement-disjointness, derivation-as-term, and
involution facts that already exist in the ZD graph.  Each is
realized as a $\sim 15$-line compiler diff, a stdlib module, a
Lean theorem, a Sounio benchmark, and a compile-fail test.  The
resulting hexad is algebraically closed and operationally usable.

\begin{thebibliography}{99}

\bibitem{ilharco2023editing}
G.~Ilharco, M.~T. Ribeiro, M.~Wortsman, S.~Gururangan, L.~Schmidt, H.~Hajishirzi,
  and A.~Farhadi.
\newblock Editing models with task arithmetic.
\newblock In \emph{ICLR}, 2023.

\bibitem{wortsman2022model}
M.~Wortsman et al.
\newblock Model soups: averaging weights of multiple fine-tuned models improves
  accuracy without increasing inference time.
\newblock In \emph{ICML}, 2022.

\bibitem{yadav2023ties}
P.~Yadav, D.~Tam, L.~Choshen, C.~Raffel, and M.~Bansal.
\newblock TIES-merging: resolving interference when merging models.
\newblock In \emph{NeurIPS}, 2023.

\bibitem{necula1997proofcarrying}
G.~C. Necula.
\newblock Proof-carrying code.
\newblock In \emph{POPL}, 1997.

\bibitem{leroy2009compcert}
X.~Leroy.
\newblock Formal verification of a realistic compiler.
\newblock \emph{CACM}, 52(7), 2009.

\bibitem{kumar2014cakeml}
R.~Kumar, M.~O. Myreen, M.~Norrish, and S.~Owens.
\newblock CakeML: a verified implementation of ML.
\newblock In \emph{POPL}, 2014.

\bibitem{seshia2022toward}
S.~A. Seshia, D.~Sadigh, and S.~S. Sastry.
\newblock Toward verified artificial intelligence.
\newblock \emph{CACM}, 65(7), 2022.

\bibitem{elhage2021mathematical}
N.~Elhage et al.
\newblock A mathematical framework for transformer circuits.
\newblock \emph{Transformer Circuits Thread}, 2021.

\bibitem{mitchell2022memorybased}
E.~Mitchell, C.~Lin, A.~Bosselut, C.~D. Manning, and C.~Finn.
\newblock Memory-based model editing at scale.
\newblock In \emph{ICML}, 2022.

\bibitem{katsumata2014parametric}
S.~Katsumata.
\newblock Parametric effect monads and semantics of effect systems.
\newblock In \emph{POPL}, 2014.

\end{thebibliography}

\end{document}
